% Job posting: https://ca.linkedin.com/jobs/view/computer-vision-engineering-at-geotab-4465773767
\documentclass[11pt]{article}
\usepackage[left=1in,top=1in,right=1in,bottom=1in]{geometry}
\usepackage[hidelinks]{hyperref}
\usepackage{setspace}
\usepackage[T1]{fontenc}
\usepackage{newtxtext,newtxmath}
\ifdefined\pdfgentounicode
  \input{glyphtounicode}
  \pdfgentounicode=1
\fi
\setstretch{1.1}
\pagenumbering{gobble}
\begin{document}
\pagestyle{empty}

\noindent
Nishal Stanislaus Silva, PhD \\
Toronto, ON \\
nishal.silva@hotmail.com \,|\, +1 613 200 2030 \\
\href{https://nishal.xyz}{nishal.xyz} \,|\, \href{https://www.linkedin.com/in/nishal-silva}{linkedin.com/in/nishal-silva} \,|\, \href{https://github.com/ns2max}{github.com/ns2max}

\vspace{1em}

\noindent Dear Hiring Manager,

\vspace{1em}

\noindent For 5.5 years at MAS Holdings, South Asia's largest apparel manufacturer, I built and ran real-time ingestion and ETL for high-frequency IoT and operational time-series data across manufacturing sites in South Asia, North America and Africa -- multi-site data pipelines processing continuous sensor streams in production. Alongside that, I designed and deployed production computer-vision inspection systems: an automated multilingual care-label QC line using OpenCV template and feature matching with an explainable defect overlay driving a conveyor and PLC reject mechanism (99.5\% reduction in inspection time), and a loom-side fabric structural-defect system using a microscopic camera array (50\% reduction in operator headcount). That combination -- high-frequency, multi-site IoT data engineering paired with embedded, production-grade computer vision -- is the same core discipline that fleet telematics and driver-safety analytics run on, and it is what drew me to this role at Geotab.

\vspace{1em}

\noindent My PhD and postdoctoral work at the University of Trento extended that discipline into real-time embedded inference under hard constraints. I own the full pipeline from acquisition and DSP/feature engineering through training, evaluation and edge deployment, and I published a detector running in 14 ms on a Raspberry Pi 4 (F1 0.76, 31.4\% CPU utilization, 74x faster than a DTW baseline). That work was about making models fast, accurate and power-efficient enough to run continuously on constrained embedded hardware -- the same set of tradeoffs that in-vehicle telematics devices demand. My C++17 and Python inference engines from MAS Holdings, deployed on embedded hardware on the factory floor, add production systems engineering to that research background.

\vspace{1em}

\noindent I have not worked in the automotive or fleet-telematics industry specifically. What I bring instead is direct, hands-on experience with the two disciplines the role is built on -- multi-site industrial IoT data pipelines at scale and embedded, production computer vision -- and I am confident that experience transfers directly to dashcam and telematics-device analytics.

\vspace{1em}

\noindent I am a Canadian Permanent Resident, eligible to work in Canada without sponsorship, based in Toronto and within easy reach of Geotab's Oakville office, and available to start immediately. I would welcome the opportunity to discuss how my background can contribute to your computer vision team.

\vspace{1em}

\noindent Sincerely, \\
Nishal Stanislaus Silva

\end{document}
